\documentclass[12pt,a4paper]{article}
\usepackage[a4paper,margin=18mm]{geometry}
\usepackage[T2A]{fontenc}
\usepackage[utf8]{inputenc}
\usepackage[russian]{babel}
\usepackage{amsmath,amssymb,mathtools}
\usepackage{array,booktabs,longtable}
\usepackage{xcolor}
\usepackage{hyperref}
\usepackage{tikz}

\hypersetup{
  colorlinks=true,
  linkcolor=blue!55!black,
  urlcolor=blue!55!black
}
\setlength{\parindent}{0pt}
\setlength{\parskip}{0.55em}
\setlength{\tabcolsep}{3pt}
\renewcommand{\arraystretch}{1.25}

\begin{document}

\thispagestyle{empty}
\begin{center}
\vspace*{0.18\textheight}
{\LARGE\bfseries Ревью домашней работы\par}
\vspace{2.2cm}
{\large Автор: Лёвин Артём Александрович\par}
\vspace{0.8cm}
{\large 7 сентября 2026 года\par}
\vfill
\begin{tikzpicture}[x=1cm,y=1cm]
  \draw[blue!45!black,line width=0.8pt]
    (-4,0) .. controls (-2.6,0.65) and (-1.3,-0.65) .. (0,0)
           .. controls (1.3,0.65) and (2.6,-0.65) .. (4,0);
\end{tikzpicture}
\vspace{1.0cm}
\end{center}

\newpage
\section*{Сводная таблица}
\small
\begin{center}
\begin{tabular}{@{}>{\centering\arraybackslash}p{14mm} p{29mm} p{49mm} p{30mm} p{30mm}@{}}
\toprule
\textbf{№ задания} & \textbf{Тема} & \textbf{Суть ошибки} & \textbf{Ваш ответ} & \textbf{Правильный ответ} \\
\midrule
\hyperref[task:1]{1} & Действия с дробями; периодическая десятичная дробь & Вычисление доведено до правильной обыкновенной дроби, но требуемая десятичная запись отсутствует. & $\frac{11}{45}$ & $0{,}2(4)$ \\
\addlinespace
\hyperref[task:2]{2} & Деление дробей; десятичная запись & Получена правильная дробь $\frac45$, но ответ не записан десятичной дробью. & $\frac45$ & $0{,}8$ \\
\addlinespace
\hyperref[task:3]{3} & Смешанные, неправильные и десятичные дроби & Неправильная дробь получена верно, но $\frac{17}{5}$ неверно переведена в десятичную дробь. & $\frac{17}{5}$; $7{,}7$ & $\frac{17}{5}$; $3{,}4$ \\
\addlinespace
\hyperref[task:4-1]{4(1)} & Сложение десятичных дробей столбиком & При переписывании первого слагаемого $2{,}34$ заменено на $3{,}34$, поэтому посчитан другой пример. & $6{,}49$ & $5{,}49$ \\
\bottomrule
\end{tabular}
\end{center}
\normalsize

\newpage
\vspace*{\fill}
\begin{center}
{\LARGE\bfseries Разбор ошибок}
\end{center}
\vspace*{\fill}

\newpage
\phantomsection\label{task:1}
\section*{Задание 1}

\textbf{Текст задания.}
Вычислите и запишите ответ десятичной дробью:
\[
\frac{7}{15}-\frac{2}{9}.
\]

\textbf{Ваше решение.}
В записи получено:
\[
\frac{7}{15}-\frac{2}{9}
=\frac{63-30}{135}
=\frac{33}{135}
=\frac{11}{45}.
\]

\textbf{Ваш ответ.}
\[
\frac{11}{45}.
\]

\textcolor{red}{\textbf{Ошибка:}} не выполнено требование о десятичной форме ответа.

\textbf{Где именно возникает ошибка.}
Последний правильный шаг --- получение дроби
\[
\frac{11}{45}.
\]
Первый неполный шаг --- на этом вычисление заканчивается, хотя по условию необходимо записать результат десятичной дробью.

\textbf{Почему это неверно.}
По значению дробь $\frac{11}{45}$ найдена верно. Ошибка относится не к вычитанию, а к форме окончательного ответа. Дробь $\frac{11}{45}$ нельзя превратить в конечную десятичную дробь: после сокращения в её знаменателе остаётся множитель $3$, поэтому десятичная запись будет бесконечной периодической. При делении $11$ на $45$ сначала получается цифра $2$ после запятой, а затем остаток $20$ повторяется, поэтому цифра $4$ дальше повторяется бесконечно.

\textcolor{blue!60!black}{\textbf{Ключевая идея:}} если условие требует десятичную дробь, после получения обыкновенной дроби нужно проверить её десятичную запись; она может быть не только конечной, но и периодической.

\textbf{Исправление от последнего правильного шага.}
Продолжаем с уже верного результата:
\[
\frac{11}{45}=11:45=0{,}24444\ldots=0{,}2(4).
\]

\textbf{Полное правильное решение.}
Наименьший общий знаменатель чисел $15$ и $9$ равен $45$:
\[
\frac{7}{15}=\frac{21}{45},
\qquad
\frac{2}{9}=\frac{10}{45}.
\]
Поэтому
\[
\frac{7}{15}-\frac{2}{9}
=\frac{21}{45}-\frac{10}{45}
=\frac{11}{45}.
\]
Теперь делим $11$ на $45$. Целая часть равна нулю. Для десятых берём $110$: $110:45=2$, остаток $20$. Для сотых берём $200$: $200:45=4$, остаток снова $20$. Такой же остаток будет появляться дальше, поэтому цифра $4$ повторяется:
\[
\frac{11}{45}=0{,}24444\ldots=0{,}2(4).
\]

\textbf{Проверка.}
Пусть $x=0{,}24444\ldots$. Тогда
\[
10x=2{,}4444\ldots,\qquad 100x=24{,}4444\ldots.
\]
Вычитаем:
\[
100x-10x=22,
\qquad
90x=22,
\qquad
x=\frac{22}{90}=\frac{11}{45}.
\]
Получена исходная дробь, значит десятичная запись верна.

\textcolor{teal!60!black}{\textbf{Итог:}}
\[
\boxed{0{,}2(4)}.
\]

\textbf{Задача на закрепление.}
Вычислите $\frac{4}{15}-\frac{1}{9}$ и запишите ответ десятичной дробью.

\newpage
\phantomsection\label{task:2}
\section*{Задание 2}

\textbf{Текст задания.}
Выполните деление и запишите ответ десятичной дробью:
\[
\frac{14}{25}:\frac{7}{10}.
\]

\textbf{Ваше решение.}
В записи деление правильно заменено умножением на обратную дробь и получено:
\[
\frac{14}{25}:\frac{7}{10}
=\frac{14}{25}\cdot\frac{10}{7}
=\frac45.
\]

\textbf{Ваш ответ.}
\[
\frac45.
\]

\textcolor{red}{\textbf{Ошибка:}} ответ не переведён в требуемую десятичную форму.

\textbf{Где именно возникает ошибка.}
Последний правильный шаг:
\[
\frac45.
\]
Первый неполный шаг --- завершение решения на обыкновенной дроби, несмотря на прямое требование условия записать ответ десятичной дробью.

\textbf{Почему это неверно.}
Само деление выполнено верно: чтобы разделить на дробь $\frac{7}{10}$, нужно умножить на обратную дробь $\frac{10}{7}$. Однако дробь $\frac45$ не является требуемой формой ответа. Здесь перевод особенно прост: знаменатель $5$ можно превратить в $10$, умножив числитель и знаменатель на $2$.

\textcolor{blue!60!black}{\textbf{Ключевая идея:}} после правильного действия с обыкновенными дробями нужно отдельно выполнить последнее требование условия --- представить результат десятичной дробью.

\textbf{Исправление от последнего правильного шага.}
\[
\frac45=\frac{8}{10}=0{,}8.
\]

\textbf{Полное правильное решение.}
Деление на дробь заменяем умножением на обратную дробь:
\[
\frac{14}{25}:\frac{7}{10}
=\frac{14}{25}\cdot\frac{10}{7}.
\]
Сокращаем $14$ и $7$:
\[
\frac{14}{7}=2.
\]
Получаем
\[
\frac{2\cdot 10}{25}=\frac{20}{25}=\frac45.
\]
Переводим результат в десятичную дробь:
\[
\frac45=\frac8{10}=0{,}8.
\]

\textbf{Проверка.}
Если частное равно $0{,}8$, то произведение частного на делитель должно дать делимое:
\[
0{,}8\cdot\frac{7}{10}=0{,}56=\frac{56}{100}=\frac{14}{25}.
\]
Проверка совпала.

\textcolor{teal!60!black}{\textbf{Итог:}}
\[
\boxed{0{,}8}.
\]

\textbf{Задача на закрепление.}
Вычислите $\frac{9}{20}:\frac{3}{8}$ и запишите ответ десятичной дробью.

\newpage
\phantomsection\label{task:3}
\section*{Задание 3}

\textbf{Текст задания.}
Переведите число
\[
3\frac25
\]
в следующие виды: 1) неправильная дробь; 2) десятичная дробь.

\textbf{Ваше решение.}
Неправильная дробь получена верно:
\[
3\frac25=\frac{17}{5}.
\]
После этого для десятичной формы записан результат $7{,}7$.

\textbf{Ваш ответ.}
\[
\frac{17}{5};\qquad 7{,}7.
\]

\textcolor{red}{\textbf{Ошибка:}} неверный перевод дроби $\frac{17}{5}$ в десятичную дробь.

\textbf{Где именно возникает ошибка.}
Последний правильный шаг:
\[
3\frac25=\frac{17}{5}.
\]
Первый неверный шаг --- получение из $\frac{17}{5}$ числа $7{,}7$. Эти числа не равны.

\textbf{Почему это неверно.}
Чтобы получить неправильную дробь из смешанного числа, действительно нужно умножить целую часть $3$ на знаменатель $5$ и прибавить числитель $2$:
\[
3\cdot5+2=17.
\]
Эта часть выполнена правильно. Далее дробь $\frac{17}{5}$ удобно привести к знаменателю $10$. Для этого нужно умножить \emph{и числитель, и знаменатель} на одно и то же число $2$; только тогда значение дроби не меняется:
\[
\frac{17}{5}=\frac{17\cdot2}{5\cdot2}=\frac{34}{10}.
\]
Дробь $\frac{34}{10}$ равна $3{,}4$, а не $7{,}7$.

\textcolor{blue!60!black}{\textbf{Ключевая идея:}} при приведении дроби к знаменателю $10$, $100$, $1000$ числитель и знаменатель умножаются на один и тот же множитель.

\textbf{Исправление от последнего правильного шага.}
Начинаем с верной дроби $\frac{17}{5}$:
\[
\frac{17}{5}=\frac{17\cdot2}{5\cdot2}=\frac{34}{10}=3{,}4.
\]

\textbf{Полное правильное решение.}
Для неправильной дроби:
\[
3\frac25=\frac{3\cdot5+2}{5}=\frac{17}{5}.
\]
Для десятичной дроби:
\[
\frac{17}{5}=\frac{34}{10}=3{,}4.
\]
Следовательно, два требуемых вида имеют вид
\[
\frac{17}{5}\qquad\text{и}\qquad 3{,}4.
\]

\textbf{Проверка.}
\[
3{,}4=3+0{,}4=3+\frac{4}{10}=3+\frac25=3\frac25.
\]
Получили исходное смешанное число.

\textcolor{teal!60!black}{\textbf{Итог:}}
\[
\boxed{\frac{17}{5};\ 3{,}4}.
\]

\textbf{Задача на закрепление.}
Переведите число $4\frac35$ в неправильную дробь и в десятичную дробь.

\newpage
\phantomsection\label{task:4-1}
\section*{Задание 4(1)}

\textbf{Текст задания.}
Выполните действие столбиком:
\[
2{,}34+3{,}15.
\]

\textbf{Ваше решение.}
В записи пример переписан как
\[
3{,}34+3{,}15,
\]
после чего для него получено $6{,}49$.

\textbf{Ваш ответ.}
\[
6{,}49.
\]

\textcolor{red}{\textbf{Ошибка:}} при переписывании изменено первое слагаемое.

\textbf{Где именно возникает ошибка.}
Последний корректный элемент записи --- второе слагаемое $3{,}15$ и знак сложения перенесены верно. Первый неверный шаг --- число $2{,}34$ из условия заменено на $3{,}34$. После этого арифметика относится уже к другому примеру.

\textbf{Почему это неверно.}
В столбик можно безошибочно сложить только те числа, которые даны в условии. Изменение одной цифры в слагаемом меняет значение всей суммы на единицу. Поэтому результат $6{,}49$ верен для выражения $3{,}34+3{,}15$, но не для исходного выражения $2{,}34+3{,}15$. Перед вычислением необходимо сверить переписанные числа с условием и поставить запятые строго друг под другом.

\textcolor{blue!60!black}{\textbf{Ключевая идея:}} при вычислении столбиком сначала точно переносим исходные числа, затем выравниваем разряды: сотые под сотыми, десятые под десятыми, единицы под единицами.

\textbf{Исправление.}
Записываем исходный пример без изменения первого слагаемого:
\[
\begin{array}{r}
  2{,}34\\
+ 3{,}15\\
\hline
  5{,}49
\end{array}
\]

\textbf{Полное правильное решение.}
Складываем по разрядам справа налево. Сотые:
\[
4+5=9.
\]
Десятые:
\[
3+1=4.
\]
Единицы:
\[
2+3=5.
\]
Запятая в ответе располагается под запятыми слагаемых. Получаем
\[
2{,}34+3{,}15=5{,}49.
\]

\textbf{Проверка.}
Вычтем второе слагаемое из найденной суммы:
\[
5{,}49-3{,}15=2{,}34.
\]
Получено первое слагаемое, значит сумма найдена верно.

\textcolor{teal!60!black}{\textbf{Итог:}}
\[
\boxed{5{,}49}.
\]

\textbf{Задача на закрепление.}
Выполните столбиком: $4{,}27+2{,}58$.

\end{document}
